% titel: Blatt 5, A4: Periodizität im PSE (Trends)
% skript: Mortimer Kap. 7 (chemie: Periodensystem, Ionisierungsenergie)

Atomradius, Ionisierungsenergie, Elektronenaffinität und Metallcharakter hängen
systematisch von der Position im PSE ab --- getrieben durch Kernladung $Z$,
Schalenzahl und Abschirmung.

\begin{werkzeugkasten}[Die vier Trends]
\begin{itemize}
  \item \textbf{In einer Periode} (links $\to$ rechts): Kernladung nimmt zu,
        Schalenzahl bleibt gleich $\to$ Atomradius nimmt ab; Bindung der
        Valenzelektronen wird fester $\to$ Ionisierungsenergie nimmt zu;
        Elektronenaffinität nimmt zu; Metallcharakter nimmt ab.
  \item \textbf{In einer Gruppe} (oben $\to$ unten): Schalenzahl und damit
        Atomradius nehmen zu, Abschirmung wächst $\to$ Ionisierungsenergie
        nimmt ab; Elektronenaffinität nimmt ab; Metallcharakter nimmt zu.
  \item \textbf{Übergangsreihen} (3d, 4d, 5d): Die zusätzlichen Elektronen
        landen in $(n-1)d$ und schirmen die steigende Kernladung nur
        unvollständig, aber konstant ab $\to$ Radius und Ionisierungsenergie
        ändern sich nur schwach über die Reihe.
  \item \textbf{Ausnahmen} bei halb-/vollbesetzten Unterschalen: besonders
        stabile Konfigurationen verursachen Maxima/Minima (z.\ B. bei Be und N).
\end{itemize}
\textit{Skript:} Mortimer Kap. 7 (Periodizität).
\end{werkzeugkasten}

\begin{achtung}[Darauf achten]
\begin{itemize}
  \item Effektive Kernladung $Z_{\mathrm{eff}}$ ist der Schlüssel: sie wächst
        in der Periode, bleibt in der Gruppe annähernd gleich.
  \item RE/Ra-Trends: Radius und Ionisierungsenergie hängen umgekehrt
        zusammen; nicht verwechseln.
  \item Elektronenaffinität wird oft als \textit{freigesetzte} Energie (\si{\kilo\joule\per\mol})
        angegeben --- im Vorzeichen achten.
  \item Bei Be ($2s^2$ voll) und N ($2p^3$ halb) liegt die 1.\ Ionisierungs-
        energie höher als der Nachbartrend erwarten lässt.
\end{itemize}
\end{achtung}
